\section{模型评价、推广与结论}

\subsection{模型优点}

\begin{enumerate}
  \item \textbf{候选与证明分离。} 四种 Q1 算法共享独立验证器；Q2/Q3 使用完整时间单元分类，不把采样点或优化器成功标志当作连续覆盖证明。
  \item \textbf{证据分层明确。} 覆盖下界、覆盖上界、总暴露、最大连续暴露和未解析区间分别序列化，联合增益使用最佳单烟幕基线计算。
  \item \textbf{信息因果性可审计。} Q4 只在揭示、认证和容量条件同时满足时调度，hindsight 被明确限制为离线上界。
  \item \textbf{结果可复现。} CSV、JSON、PNG 和检查脚本使用同一个冻结基线标识，论文表格可以回溯到仓库文件。
\end{enumerate}

\subsection{模型局限}

第一，Q1 的四算法比较使用固定预算和固定场景矩阵，不能替代对连续可行域的全局求解证明。第二，Q2 的事件候选和 SLSQP 精修是有界搜索，候选集之外的事件组合没有被声称覆盖。第三，Q3 的主解释固定为三架 UAV 每架一枚，路径冲突的安全距离参数按代码为零，因此通过证书的现实含义受限。第四，Q4 任务包规模小，三种策略相同的数值不能外推到更大的动态任务流。第五，灵敏度部分包含固定路径、固定候选和实验反事实，不能替代完整概率鲁棒性分析。

\subsection{推广方向}

若需要扩大模型，可在保持“候选—验证—证书”边界的前提下，引入更密的事件候选、区间算术或可证明的全局搜索；若需要更严格的多机安全约束，应在路径证书中设置非零安全距离并重新生成 Q3 工件；若需要研究更大规模 Q4，可把认证任务包接入明确的在线整数规划或近似调度器，但必须重新定义可用信息和上界含义。上述方向都是后续工作，不属于本 Stage 8 已提交的科学结果。

\subsection{结论}

本文完成了一条可审计的多阶段建模链：Q1 给出四算法 benchmark 和独立 verifier；Q2 用有限事件候选、SLSQP 精修与保守联合覆盖认证；Q3 在“三架 UAV、每架一枚”的解释下检查路径冲突并进行带 $\varepsilon$ 的词典序筛选；Q4 在认证任务包上进行因果 rolling、因果 greedy 与离线 hindsight 上界比较。结果文件支持联合覆盖、暴露区间、路径证书和调度值等定量结论，但不支持完整 SIP、完整 Pareto、滚动 MILP 或连续全局最优的表述。

\subsection{Stage 8 一致性声明}

Stage 8 只新建论文、图表副本、审计和支撑材料，没有修改 Q1--Q4 算法或科学结果。若读者发现 CSV、JSON、图表与正文之间的差异，应以冻结结果文件和独立检查脚本为优先证据，并在审计表中登记，而不是在论文中自行修正数值。
